\documentclass[11pt]{article}
\newcommand{\manualtitle}{BDX Slow Control Quickstart Guide}
\newcommand{\manualshorttitle}{Quickstart Guide}
\usepackage[T1]{fontenc}
\usepackage[utf8]{inputenc}
\usepackage{lmodern}
\usepackage[a4paper,margin=24mm,headheight=15pt]{geometry}
\usepackage{microtype}
\usepackage{parskip}
\usepackage{enumitem}
\usepackage{booktabs}
\usepackage{tabularx}
\usepackage{longtable}
\usepackage{array}
\usepackage{xcolor}
\usepackage{listings}
\usepackage[most]{tcolorbox}
\usepackage{fancyhdr}
\usepackage{lastpage}
\usepackage[hidelinks]{hyperref}
\usepackage{bookmark}
\usepackage{url}

\definecolor{BDXBlue}{HTML}{1F4E79}
\definecolor{BDXLightBlue}{HTML}{EAF2F8}
\definecolor{BDXGray}{HTML}{F2F3F4}
\definecolor{BDXDarkGray}{HTML}{333333}
\definecolor{BDXOrange}{HTML}{A64B00}
\definecolor{BDXLightOrange}{HTML}{FFF2E5}
\definecolor{BDXRed}{HTML}{8B1A1A}
\definecolor{BDXLightRed}{HTML}{FCE8E8}
\definecolor{BDXGreen}{HTML}{236B2E}
\definecolor{BDXLightGreen}{HTML}{E8F5E9}

\hypersetup{
  colorlinks=true,
  linkcolor=BDXBlue,
  urlcolor=BDXBlue,
  citecolor=BDXBlue,
  pdftitle={\manualtitle},
  pdfauthor={BDX Collaboration}
}

\setlist[itemize]{leftmargin=*,topsep=3pt,itemsep=2pt}
\setlist[enumerate]{leftmargin=*,topsep=3pt,itemsep=3pt}
\renewcommand{\arraystretch}{1.15}
\setlength{\emergencystretch}{2em}

\pagestyle{fancy}
\fancyhf{}
\lhead{BDX Slow Control}
\rhead{\manualshorttitle}
\cfoot{Page \thepage\ of \pageref*{LastPage}}
\renewcommand{\headrulewidth}{0.4pt}

\lstdefinestyle{bdxbase}{
  basicstyle=\ttfamily\small,
  columns=fullflexible,
  keepspaces=true,
  frame=single,
  rulecolor=\color{black!20},
  backgroundcolor=\color{BDXGray},
  breaklines=true,
  breakatwhitespace=false,
  showstringspaces=false,
  tabsize=2,
  xleftmargin=0.5em,
  xrightmargin=0.5em,
  framexleftmargin=0.4em,
  framexrightmargin=0.4em
}
\lstdefinestyle{shell}{style=bdxbase,language=bash}
\lstdefinestyle{python}{style=bdxbase,language=Python}
\lstdefinestyle{json}{style=bdxbase,language={},morestring=[b]",stringstyle=\color{BDXGreen},morecomment=[l]{//}}
\lstset{style=bdxbase}

\newtcolorbox{importantbox}[1][Important]{
  enhanced,
  breakable,
  colback=BDXLightOrange,
  colframe=BDXOrange,
  title=\textbf{#1},
  fonttitle=\bfseries,
  boxrule=0.7pt,
  arc=1mm
}
\newtcolorbox{warningbox}[1][Warning]{
  enhanced,
  breakable,
  colback=BDXLightRed,
  colframe=BDXRed,
  title=\textbf{#1},
  fonttitle=\bfseries,
  boxrule=0.7pt,
  arc=1mm
}
\newtcolorbox{notebox}[1][Note]{
  enhanced,
  breakable,
  colback=BDXLightBlue,
  colframe=BDXBlue,
  title=\textbf{#1},
  fonttitle=\bfseries,
  boxrule=0.7pt,
  arc=1mm
}
\newtcolorbox{successbox}[1][Expected result]{
  enhanced,
  breakable,
  colback=BDXLightGreen,
  colframe=BDXGreen,
  title=\textbf{#1},
  fonttitle=\bfseries,
  boxrule=0.7pt,
  arc=1mm
}

\newcommand{\repo}{\href{https://github.com/vallarinosimone1992/bdx-slow-control}{\texttt{vallarinosimone1992/bdx-slow-control}}}
\newcommand{\file}[1]{\path{#1}}
\newcommand{\command}[1]{\texttt{\detokenize{#1}}}
\newcommand{\pv}[1]{\texttt{\detokenize{#1}}}
\newcommand{\code}[1]{\texttt{\detokenize{#1}}}
\newcommand{\documentstatus}{Development documentation for the \texttt{devel} branch}

\newcommand{\makemanualtitle}{
  \begin{titlepage}
    \centering
    \vspace*{24mm}
    {\sffamily\bfseries\color{BDXBlue}\Huge \manualtitle\par}
    \vspace{8mm}
    {\sffamily\Large BDX Slow Control\par}
    \vspace{18mm}
    \begin{tcolorbox}[width=0.82\textwidth,colback=BDXLightBlue,colframe=BDXBlue,arc=1mm]
      \centering
      \textbf{Document status:} \documentstatus\\[2mm]
      \textbf{Repository:} \repo\\[2mm]
      \textbf{Revision date:} \today
    \end{tcolorbox}
    \vfill
    {\large BDX Collaboration\par}
  \end{titlepage}
}


\begin{document}
\hypersetup{pageanchor=false}
\makemanualtitle
\hypersetup{pageanchor=true}
\pagenumbering{roman}
\tableofcontents
\clearpage
\pagenumbering{arabic}

\section{Purpose and audience}

This guide is for operators who need to start, monitor, and stop the BDX slow-control prototype without modifying its software. It covers the normal Ubuntu desktop workflow, the Phoebus operator displays, the low-voltage power supplies, the chiller, the Raspberry Pi temperature IOC, and basic fault checks.

The software uses EPICS Channel Access through caproto. The main Ubuntu host runs the aggregated IOC for the global, low-voltage power-supply, and chiller subsystems. A Raspberry Pi runs the environmental temperature IOC. Phoebus provides the graphical operator interface, and EPICS Archiver Appliance stores selected process variables for later retrieval.

\begin{warningbox}[Operational safety]
The slow-control software is not a substitute for hardwired protection, approved operating procedures, or direct verification at the equipment. Before enabling an output or changing a setpoint, confirm the intended channel, acceptable limits, load state, and detector condition.
\end{warningbox}

\section{System summary}

The current prototype uses the following fixed slow-control network endpoints.

\begin{center}
\begin{tabularx}{\textwidth}{@{}l l X@{}}
\toprule
Component & Address & Function \\
\midrule
Main IOC host & \code{172.22.50.2} & Global state, LV1, LV2, and chiller IOC groups \\
Raspberry Pi & \code{172.22.50.10} & MCP9808 environmental temperature IOC \\
LV1 & \code{172.22.50.20:9221} & TTi CPX400DP low-voltage supply \\
LV2 & \code{172.22.50.21:9221} & TTi CPX400DP low-voltage supply \\
Chiller & \code{172.22.50.60:54321} & LAUDA ECO Silver RE 1225 S \\
Archiver & \code{127.0.0.1:17665--17668} & Local Archiver Appliance services \\
\bottomrule
\end{tabularx}
\end{center}

The default operational configuration is stored in \file{config/profiles/default}. It includes the global, PSU, and chiller subsystems. The Raspberry configuration is stored separately in \file{config/profiles/raspberry}.

\section{Before starting}

\subsection{Open a local desktop terminal}

Run the normal start command from a terminal opened inside the Ubuntu graphical desktop session. The launcher needs a working \code{DISPLAY} or \code{WAYLAND_DISPLAY} and opens separate terminal windows for the IOC and for the Archiver/Phoebus workflow.

A plain SSH session without graphical desktop access is not the normal operator workflow.

\subsection{Enter the repository}

\begin{lstlisting}[style=shell]
cd ~/SlowControl/app/bdx-slow-control
\end{lstlisting}

If the repository is installed elsewhere, enter that directory or set \code{BDX_SLOW_CONTROL_ROOT} to the checkout path.

\subsection{Check the runtime address}

The untracked file \file{config/runtime.env} must identify the main IOC host. Create it once from the template if needed:

\begin{lstlisting}[style=shell]
cp config/runtime.env.example config/runtime.env
\end{lstlisting}

For the current prototype, the file should contain:

\begin{lstlisting}[style=shell]
BDX_MAIN_HOST=172.22.50.2
\end{lstlisting}

Do not use \code{127.0.0.1} for normal operation. Remote Channel Access clients cannot reach a loopback-only IOC.

\subsection{Optional connectivity checks}

Read-only network checks can be used before startup:

\begin{lstlisting}[style=shell]
ping -c 2 172.22.50.20
nc -vz 172.22.50.20 9221

ping -c 2 172.22.50.21
nc -vz 172.22.50.21 9221

ping -c 2 172.22.50.60
nc -vz 172.22.50.60 54321
\end{lstlisting}

These checks do not write EPICS PVs and do not change equipment settings.

\section{Starting the slow control}

\subsection{Start or verify the Raspberry environment IOC}

The main startup command checks whether \pv{BDX:ENV:TEMP:T00:VALUE} is reachable. If the Raspberry IOC is already running, no action is required. To start it through SSH:

\begin{lstlisting}[style=shell]
start-bdx-raspberry-ioc
\end{lstlisting}

To force a restart:

\begin{lstlisting}[style=shell]
start-bdx-raspberry-ioc --restart
\end{lstlisting}

The command runs the remote systemd action and then verifies the readiness PV over Channel Access.

\subsection{Start the main operator stack}

From the repository root:

\begin{lstlisting}[style=shell]
bdx_slow_control_start overview
\end{lstlisting}

The command performs the following actions:

\begin{enumerate}
  \item checks the Raspberry environment IOC and prints a warning if it is unavailable;
  \item opens a terminal for the main aggregated IOC unless one is already listening on port 5064;
  \item waits for the main IOC readiness PV;
  \item checks or starts the local Archiver Appliance;
  \item checks representative PSU, chiller, and environment archive connections;
  \item launches Phoebus on the requested display.
\end{enumerate}

The startup command does not enable PSU outputs, start or stop the chiller, or write hardware setpoints.

\begin{successbox}
A normal startup leaves two terminal windows open: \textbf{BDX Main IOC} and \textbf{BDX Archiver and Phoebus}. Phoebus opens the overview display. Keep the terminals available for diagnostics.
\end{successbox}

A different initial display can be selected:

\begin{lstlisting}[style=shell]
bdx_slow_control_start psu
bdx_slow_control_start chiller
bdx_slow_control_start environment
bdx_slow_control_start trends
\end{lstlisting}

\section{Using Phoebus}

\subsection{Main displays}

\begin{tabularx}{\textwidth}{@{}l X@{}}
\toprule
Display & Purpose \\
\midrule
\file{overview.bob} & Overall status and navigation \\
\file{psu.bob} & Routine LV1/LV2 operation with staged requests and live plots \\
\file{psu_expert.bob} & Complete PSU PV table and direct compatibility controls \\
\file{chiller.bob} & Routine chiller operation and live temperature plots \\
\file{chiller_expert.bob} & Complete chiller PV table and expert controls \\
\file{environment.bob} & Environmental sensor status and values \\
\file{global.bob} & Global state, update timing, interlock test, reset, and all-off \\
\file{trends.bob} & Principal live and archived trends \\
\file{all_pvs.bob} & Complete configured PV catalog \\
\bottomrule
\end{tabularx}

Operator pages intentionally expose the normal safe workflow. Expert pages contain the complete writable PV surface and should be used only when the direct behavior is understood.

\subsection{Status fields to check}

Before issuing a command, verify the relevant IOC and device status:

\begin{itemize}
  \item \pv{IOC_STATE} should be \code{RUNNING};
  \item \pv{COMM_OK} should be true;
  \item \pv{COMM_STATUS} should be \code{OK} for hardware or \code{SIMULATION} for a simulator;
  \item \pv{LAST_UPDATE} should continue changing;
  \item \pv{ERROR_CODE} should be zero and \pv{ERROR_MESSAGE} should be empty.
\end{itemize}

If communication is not healthy, do not assume that an old displayed value represents the current equipment state.

\section{Operating the low-voltage supplies}

The operator display uses a staged request/apply workflow. Editing a request field does not immediately write the instrument.

\subsection{Set voltage and current limit}

For the intended device and channel:

\begin{enumerate}
  \item confirm the device label, channel, communication status, and current output state;
  \item enter the desired value in \pv{VOLTAGE_REQUEST};
  \item enter the desired value in \pv{CURRENT_LIMIT_REQUEST};
  \item recheck both values and press the confirmed \textbf{Apply} button;
  \item verify \pv{APPLY_STATUS} and \pv{APPLY_MESSAGE};
  \item verify the hardware readbacks:
  \begin{itemize}
    \item \pv{VOLTAGE_SET_RBV};
    \item \pv{VOLTAGE_RBV};
    \item \pv{CURRENT_LIMIT_RBV}.
  \end{itemize}
\end{enumerate}

The IOC validates voltage and current together before writing either setting. The default software envelope is 0--60 V, 0--20 A, and a maximum requested voltage-current product of 420 W. A profile may impose stricter per-device or per-channel limits.

\begin{importantbox}[Setpoint versus readback]
A writable setpoint or request indicates what was requested or accepted by the IOC. A field ending in \code{_RBV} is updated from a device query and represents the hardware readback. Always verify the readback after a command.
\end{importantbox}

\subsection{Enable or disable an output}

Use the confirmed output control for the intended channel. After the command, verify both \pv{OUTPUT_RBV} and the measured voltage/current. Do not rely only on the button state.

\subsection{All-off}

Device and global all-off commands are available in Phoebus. The global command is intended to disable managed power outputs and requires confirmation.

\begin{warningbox}[All-off scope]
Software all-off acts through the configured drivers. It does not replace hardware emergency-off systems, and its effectiveness depends on communication with the equipment. Verify the physical result when safety or detector protection depends on the output state.
\end{warningbox}

\section{Operating the chiller}

\subsection{Apply a temperature setpoint}

\begin{enumerate}
  \item verify chiller communication and confirm that the displayed temperatures are updating;
  \item enter the target in \pv{SETPOINT_REQUEST};
  \item press the confirmed setpoint apply control;
  \item verify the apply status and the applied setpoint readback;
  \item monitor controlled temperature, bath temperature, and deviation alarms.
\end{enumerate}

The current configured range is 5--40 degrees C. The default deviation thresholds are 0.2 degrees C for warning and 0.5 degrees C for alarm.

\subsection{Start or stop circulation}

The chiller does not start automatically when the IOC starts. Use the confirmed \textbf{Start} or \textbf{Stop} control, then verify the running-state readback and temperature evolution.

Pressure and external-temperature measurements are disabled in the current main-server profile because those values are not available from the present setup. Disabled optional measurements are not queried and are not shown on the operator page.

\section{Viewing archived history}

Phoebus Data Browser combines live Channel Access data with Archiver Appliance retrieval when the PV is registered and actively sampled. Operator pages provide principal live plots and \textbf{Full history} actions for archived PVs.

\begin{importantbox}[Archive limitation]
History is available only from the time a PV was successfully registered and sampled. Missing past samples cannot be reconstructed retroactively.
\end{importantbox}

If a recent value appears live but not in history, check:

\begin{itemize}
  \item that the Archiver terminal reported the relevant subsystem as ready;
  \item that the PV is registered and connected in Archiver Appliance;
  \item that the main host, Raspberry Pi, and Archiver host clocks are synchronized;
  \item that the selected plot time range includes the expected timestamp.
\end{itemize}

\section{Stopping the software}

The software components can be stopped independently:

\begin{lstlisting}[style=shell]
bdx_slow_control_kill_phoebus
bdx_slow_control_kill_archiver
bdx_slow_control_kill_ioc
\end{lstlisting}

To stop the complete local software stack:

\begin{lstlisting}[style=shell]
./scripts/kill_slow_control_all.sh
\end{lstlisting}

Normal shutdown is graceful. The \code{--force} option is reserved for processes that do not terminate normally.

\begin{warningbox}[Software shutdown does not make the hardware safe]
Stopping Phoebus, Archiver Appliance, or the IOC does not switch off PSU outputs, send an automatic all-off, stop the chiller, or change setpoints. Put the equipment in the intended state before stopping the IOC, and verify that state directly.
\end{warningbox}

\section{Basic troubleshooting}

\begin{longtable}{@{}p{0.29\textwidth}p{0.65\textwidth}@{}}
\toprule
Symptom & Check or action \\
\midrule
\endhead
Startup reports no graphical display & Run \command{bdx_slow_control_start} from a terminal in the Ubuntu desktop session. Confirm that \code{DISPLAY} or \code{WAYLAND_DISPLAY} is set. \\
\addlinespace
Runtime environment not found & Create \file{config/runtime.env} from the example and set \code{BDX_MAIN_HOST=172.22.50.2}. \\
\addlinespace
Phoebus launcher not found & Confirm the Phoebus installation. Set \code{BDX_PHOEBUS_HOME} or pass \code{--phoebus-home}. The current default is \file{~/SlowControl/css/phoebus-4.7.4-SNAPSHOT}. \\
\addlinespace
Raspberry warning at startup & Run \command{start-bdx-raspberry-ioc}. If it fails, inspect the remote service with \command{ssh -t pi@172.22.50.10 sudo journalctl -u bdx-environment-ioc -n 100 --no-pager}. \\
\addlinespace
A PSU does not communicate & Check the instrument Ethernet cable and front-panel network configuration. Run \command{ping} and \command{nc -vz <address> 9221}. Check whether another client already holds the instrument socket. \\
\addlinespace
Chiller does not communicate & Check \code{172.22.50.60:54321}, the chiller network configuration, and the IOC terminal error. \\
\addlinespace
PV is not found & Confirm the correct Channel Access address list, that the relevant IOC is running, and that no second IOC exposes the same PV names. \\
\addlinespace
Values are live but history is empty & Confirm Archiver registration, connection state, retrieval availability, and synchronized clocks. \\
\addlinespace
Display looks stale & Check \pv{LAST_UPDATE}, \pv{COMM_OK}, and the IOC terminal. A visible old value is not proof of current communication. \\
\bottomrule
\end{longtable}

\section{Command reference}

\begin{tabularx}{\textwidth}{@{}l X@{}}
\toprule
Command & Function \\
\midrule
\command{bdx_slow_control_start overview} & Start or reuse the main IOC, check Archiver, and open Phoebus \\
\command{start-bdx-raspberry-ioc} & Start the Raspberry environment IOC and verify its readiness PV \\
\command{start-bdx-raspberry-ioc --restart} & Restart the Raspberry service and verify it \\
\command{bdx_slow_control_kill_phoebus} & Stop BDX slow-control Phoebus processes \\
\command{bdx_slow_control_kill_archiver} & Stop the repository-managed local Archiver Appliance \\
\command{bdx_slow_control_kill_ioc} & Stop BDX main IOC processes \\
\command{./scripts/kill_slow_control_all.sh} & Stop all local slow-control software components \\
\command{caproto-get <PV>} & Read one PV from the command line \\
\command{caproto-put <PV> <value>} & Write one PV; use only with an understood command or setpoint \\
\bottomrule
\end{tabularx}

\section{Further documentation}

Detailed architecture, configuration, development, testing, and extension procedures are provided in the \emph{BDX Slow Control Developer Manual}. Raspberry-specific deployment details remain in \file{docs/raspberry.md}; Archiver-specific deployment details remain in \file{deploy/archiver-appliance/README.md}.

\end{document}
